\begin{conclusion}[\texorpdfstring{$(4)$}{(4)}-division 的重量像：临界集合的 \(y\)-值由不可约十二次多项式支配]\label{con:xi-terminal-zm-4division-weight-image}
在同一椭圆曲线模型
\[
E:\ Y^{2}=X^{3}-X+\frac14,
\qquad
y:=X^{2}+Y-\frac12
\]
及其诱导的谱四次平面曲线
\[
\Pi(X,y):=X^{4}-X^{3}-(2y+1)X^{2}+X+y(y+1)=0
\]
的背景下，令
\[
N_{6}(X):=2X^{6}-10X^{4}+10X^{3}-10X^{2}+2X+1
\]
（其零点给出 \(x\bigl(E[4]\setminus E[2]\bigr)\)，等价于结论 \ref{con:xi-terminal-zm-lattes-4division-certificate} 的 \(4\)-division 证书多项式 \(\psi_{4}/(2Y)\) 在 \(X\)-坐标上的表示）。
定义
\[
T_{4}(y):=\operatorname{Res}_{X}\bigl(N_{6}(X),\Pi(X,y)\bigr)\in\ZZ[y].
\]
则 \(T_{4}\) 为不可约的十二次整系数多项式，并且
\[
y\bigl(E[4]\setminus E[2]\bigr)=\{\,u\in\overline{\QQ}:\ T_{4}(u)=0\,\}.
\]
此外，\(T_{4}\) 与其判别式具有强烈的算术刚性；其展开式与判别式素分解如下：
% Auto-generated by scripts/exp_fold_zm_elliptic_4division_weight_image_audit.py
\[
T_{4}(y)=16 y^{12} - 224 y^{11} + 800 y^{10} + 2832 y^{9} - 5912 y^{8} - 46784 y^{7} - 59144 y^{6} - 9656 y^{5} + 29776 y^{4} + 27576 y^{3} + 6468 y^{2} + 312 y - 725.
\]
\[
\operatorname{Disc}\!\bigl(T_{4}\bigr)=2^{56}\cdot 37^{23}\cdot 59^{6}\cdot(652225136287)^{2}.
\]
\[
T_{4}(y)\equiv 16\,(y-6)^{10}\,(y^{2}+9y+6)\ (\mathrm{mod}\ 37).
\]

特别地，\(\operatorname{Disc}(T_{4})\) 中素数 \(37\) 的指数为 \(23\)，且其平方自由部分等于 \(37\)；从而该重量像所张成的数域在 \(37\) 处发生高阶分歧。
\end{conclusion}
\begin{proof}[证明要点]
由 \(P\in E[4]\setminus E[2]\) 当且仅当 \([2]P\in E[2]\setminus\{O\}\)，可知其 \(X\)-坐标满足 \(4\)-division 证书六次 \(N_{6}(X)=0\)。
另一方面，由 \(y=X^{2}+Y-\tfrac12\) 消去 \(Y\) 得 \(\Pi(X,y)=0\)（结论 \ref{con:xi-terminal-zm-elliptic-normalization}）。
因此 \(y(E[4]\setminus E[2])\) 正是消去 \(X\) 后的参数像，结果式 \(\operatorname{Res}_{X}(N_{6},\Pi)\) 给出其最小整系数消元证书；
不可约性与判别式分解可由脚本 \texttt{scripts/exp\_fold\_zm\_elliptic\_4division\_weight\_image\_audit.py} 复核。证毕。
\end{proof}

